\chapter{Thiết kế và xây dựng hệ thống TechStore}
\label{chap:c3_design}

Chương này trình bày quá trình phân tích yêu cầu, thiết kế kiến trúc tổng thể, cơ sở dữ liệu, sơ đồ trạng thái, các luồng xử lý chính và pipeline RAG cho chatbot AI. Trọng tâm của chương là pipeline RAG vì đây là thành phần cốt lõi và sáng tạo nhất của hệ thống.

\section{Phân tích yêu cầu}
\label{sec:c3_requirements}

Phần này xác định các tác nhân, yêu cầu chức năng và phi chức năng của hệ thống.

\subsection{Xác định tác nhân}

Hệ thống TechStore phục vụ bốn nhóm tác nhân chính:
\begin{itemize}
  \item \textbf{Khách vãng lai} (guest): duyệt sản phẩm, tìm kiếm, xem chi tiết, thêm vào giỏ hàng tạm thời (lưu trên trình duyệt), và tương tác với chatbot AI mà không cần tài khoản.
  \item \textbf{Khách hàng} (customer): có đầy đủ quyền như khách vãng lai, cộng thêm đặt hàng, thanh toán trực tuyến, theo dõi đơn hàng, viết đánh giá và quản lý thông tin cá nhân.
  \item \textbf{Nhân viên} (staff): truy cập back-office để quản lý sản phẩm, danh mục, đơn hàng, tồn kho, hoàn tiền và kiểm duyệt đánh giá.
  \item \textbf{Quản trị viên} (admin): truy cập back-office ở chế độ xem đối với dữ liệu nghiệp vụ, độc quyền quản lý tài khoản người dùng và xem báo cáo thống kê.
\end{itemize}

\subsection{Yêu cầu chức năng}

Hệ thống có 28 ca sử dụng chia thành 7 nhóm: duyệt sản phẩm (4 UC), xác thực (5 UC), giỏ hàng (3 UC), đặt hàng và thanh toán (5 UC), chatbot AI (2 UC), quản trị nghiệp vụ cho staff (6 UC), và quản trị hệ thống cho admin (3 UC). Chatbot AI là trọng tâm với yêu cầu: hiểu câu hỏi tiếng Việt kể cả viết tắt, truy xuất sản phẩm từ vector store, sinh phản hồi tự nhiên kèm sản phẩm gợi ý, và từ chối câu hỏi ngoài phạm vi. Phân quyền RBAC: staff có toàn quyền CRUD nghiệp vụ nhưng không truy cập tài khoản người dùng; admin độc quyền quản lý users và analytics.

\subsection{Yêu cầu phi chức năng}

Về hiệu năng, API phải phản hồi trong vòng 200ms với ít hơn 100 người dùng đồng thời cho các thao tác CRUD thông thường. Chatbot AI có thể mất 2 đến 5 giây do phụ thuộc vào API của nhà cung cấp LLM, điều này được chấp nhận và thông báo với người dùng thông qua trạng thái loading. Vector store JSON phải phục vụ Hybrid Search trong vòng 100ms cho catalog dưới 10.000 sản phẩm.

Về bảo mật, hệ thống phải triển khai xác thực JWT dual-token với access token 7 ngày và refresh token 30 ngày, hash mật khẩu bằng bcrypt với cost factor 12, bảo vệ chống CSRF qua cơ chế httpOnly cookie với thuộc tính \texttt{SameSite=Strict}, và mã hóa HMAC cho callback từ cổng thanh toán. Rate limiting được áp dụng theo từng nhóm endpoint: API chung giới hạn 100 request/15 phút, xác thực 10 request/60 phút, OTP 5 request/15 phút mỗi email, và chatbot 20 request/60 giây mỗi IP.

Về độ tin cậy, giao dịch đặt hàng và trừ tồn kho phải đảm bảo tính nguyên tử (ACID) để tránh overselling. Cơ chế idempotency cho IPN callback ngăn xử lý trùng lặp khi cổng thanh toán gửi lại thông báo. Chatbot phải có cơ chế dự phòng (fallback) khi LLM API không khả dụng, với chain fallback qua ba nhà cung cấp embedding để đảm bảo tính sẵn sàng cao.

Về khả năng kiểm thử, coverage tối thiểu đạt statements và branches $\geq$99,7\%, mutation score $\geq$70\% (Stryker).

\subsection{Biểu đồ ca sử dụng tổng quan}

Biểu đồ ca sử dụng được chia thành ba sơ đồ theo nhóm tác nhân (Hình~\ref{fig:usecase_overview_guest}, \ref{fig:usecase_overview_customer}, \ref{fig:usecase_overview_admin}). Khách vãng lai duyệt catalog và tương tác chatbot AI không cần tài khoản. Khách hàng có thêm đặt hàng, thanh toán và đánh giá. Back-office áp dụng kế thừa actor: staff độc quyền CRUD nghiệp vụ, admin độc quyền quản lý users.

\begin{figure}[H]
  \centering
  \fitfig[0.88]{figures/c3/usecase_overview_guest.pdf}
  \caption{Biểu đồ ca sử dụng -- Khách vãng lai (Guest)}
  \label{fig:usecase_overview_guest}
\end{figure}

\begin{figure}[H]
  \centering
  \fitfig[0.82]{figures/c3/usecase_overview_customer.pdf}
  \caption{Biểu đồ ca sử dụng -- Khách hàng (Customer)}
  \label{fig:usecase_overview_customer}
\end{figure}

\begin{figure}[H]
  \centering
  \fitfig[0.88]{figures/c3/usecase_overview_admin.pdf}
  \caption{Biểu đồ ca sử dụng -- Back-office (Nhân viên và Quản trị viên)}
  \label{fig:usecase_overview_admin}
\end{figure}

\section{Thiết kế kiến trúc hệ thống}
\label{sec:c3_architecture}

\subsection{Kiến trúc tổng quan}

Hình~\ref{fig:system_architecture} minh họa kiến trúc tổng quan của hệ thống TechStore. Hệ thống được triển khai theo kiến trúc bốn tầng.

\begin{figure}[H]
  \centering
  \fitfig[0.95]{figures/c3/system_architecture.pdf}
  \caption{Kiến trúc tổng quan hệ thống TechStore}
  \label{fig:system_architecture}
\end{figure}

Hệ thống gồm bốn tầng:
\begin{itemize}
  \item \textbf{Client}: React 19 SPA, TanStack Query + 6 Zustand stores, floating chat widget.
  \item \textbf{API Server}: Node.js\slash Express, 17 module Modular Monolith, middleware bảo mật + rate limiting.
  \item \textbf{Dữ liệu}: MySQL 8 với 25 model + vector store JSON 1024 chiều.
  \item \textbf{Dịch vụ ngoài}: LLM API, embedding chain fallback, MoMo\slash VNPay, Gmail SMTP, Google OAuth.
\end{itemize}

\subsection{Kiến trúc backend: Modular Monolith}

Hình~\ref{fig:modular_architecture} minh họa kiến trúc backend. Mỗi module là một vertical slice khép kín (Controller → Service → Repository). Các module giao tiếp qua ba cơ chế: \textbf{DI} (phụ thuộc truyền qua constructor trong \texttt{app.js})~\cite{richter2020modular}, \textbf{EventBus} (pub/sub bất đồng bộ, dùng \texttt{Promise.allSettled})~\cite{hohpe2004enterprise}, và \textbf{Shared Models} (inject model thay vì gọi API nội bộ). Lựa chọn Modular Monolith thay vì Microservices vì dự án ở giai đoạn khởi đầu với nhóm nhỏ, chưa cần scale từng module riêng lẻ.

\begin{figure}[H]
  \centering
  \fitfig{figures/c3/modular_architecture.pdf}
  \caption{Kiến trúc Modular Monolith của backend TechStore}
  \label{fig:modular_architecture}
\end{figure}

Cơ chế \textbf{UnitOfWork}~\cite{fowler2003patterns} đặc biệt quan trọng trong luồng đặt hàng. Pattern này bọc toàn bộ chuỗi thao tác (tạo order, tạo order items, trừ tồn kho) trong một transaction duy nhất với \texttt{runInTransaction}. Khi cần trừ tồn kho, \texttt{lockRow} thực thi câu lệnh \texttt{SELECT FOR UPDATE} để khóa hàng dữ liệu, ngăn hai request đồng thời cùng mua sản phẩm cuối cùng trong kho (race condition phổ biến trong hệ thống thương mại điện tử).

17 module backend chia thành hai nhóm: 12 Full DI modules (phụ thuộc truyền qua constructor, isolation tốt) và 5 Singleton modules (wrapper mỏng). Các module chính: \textit{auth}, \textit{users}, \textit{catalog}, \textit{cart}, \textit{orders}, \textit{payment}, \textit{inventory}, \textit{ai} (chatbot RAG), \textit{reviews}, \textit{wishlist}, \textit{content}, \textit{upload}, \textit{admin}, \textit{attribute}, \textit{image}, \textit{search-history}, \textit{discount-code}.

\subsection{Kiến trúc frontend: Feature-Based}

Frontend được xây dựng bằng React 19 và TypeScript theo mô hình Feature-Based với 13 feature folders, mỗi folder tự đóng gói (\texttt{api/}, \texttt{components/}, \texttt{hooks/}, \texttt{pages/}, \texttt{types/}) và không import chéo giữa các feature. TanStack Query v5 quản lý server state, 6 Zustand stores quản lý client state (auth, cart, chat, catalog, wishlist, ui). Axios singleton với hai interceptor xử lý tự động JWT injection và refresh token deduplication.

\subsection{Tổng quan API endpoints}
\label{sec:c3_api_summary}

Hệ thống cung cấp REST API chia thành 7 nhóm chính: xác thực (6 endpoints), danh mục sản phẩm (4 endpoints public), giỏ hàng (5 endpoints), đơn hàng (4 endpoints), thanh toán (5 endpoints bao gồm IPN callback), chatbot AI (3 endpoints), và quản trị back-office (10 endpoints). Mỗi endpoint được bảo vệ bởi mức xác thực phù hợp: \textit{Public}, \textit{Customer}, \textit{Staff}, \textit{Admin} hoặc \textit{Back-office} (Staff hoặc Admin). Phân quyền RBAC thực thi qua middleware \texttt{adminAuthenticate} (cho cả admin và staff vào back-office), \texttt{requireRole('staff')} và \texttt{requireSuperAdmin} phân quyền chi tiết từng route. Tài liệu API đầy đủ được sinh tự động tại \texttt{/api-docs} qua Swagger/OpenAPI 3.0.

\section{Thiết kế cơ sở dữ liệu}
\label{sec:c3_database}

\subsection{Tổng quan cấu trúc}

Cơ sở dữ liệu TechStore sử dụng MySQL 8 với charset utf8mb4 để hỗ trợ đầy đủ Unicode bao gồm tiếng Việt và emoji. Hệ thống có 25 model Sequelize tương ứng 25 bảng chính (model Image tồn tại dạng file nhưng đã loại khỏi \texttt{index.js}, không đăng ký vào ứng dụng), được quản lý qua 62 migration files theo nguyên tắc không bao giờ chỉnh sửa migration đã chạy. Tất cả bảng có trường \texttt{created\_at}, \texttt{updated\_at} và phần lớn có \texttt{deleted\_at} (paranoid/soft delete) để tránh mất dữ liệu không thể khôi phục.

\subsection{Thiết kế bảng cốt lõi}

Hình~\ref{fig:erd_core} trình bày sơ đồ ERD của 10 bảng cốt lõi.

\begin{figure}[H]
  \centering
  \fitfig{figures/c3/erd_core.pdf}
  \caption{Sơ đồ ERD các bảng cốt lõi hệ thống TechStore}
  \label{fig:erd_core}
\end{figure}

Product là thực thể trung tâm có quan hệ với hầu hết các bảng khác. Bảng \texttt{products} lưu thông tin chung của sản phẩm: tên, mô tả, thương hiệu, danh mục, ảnh đại diện, và trạng thái (active, inactive, hoặc draft). Mỗi sản phẩm có thể có nhiều biến thể trong bảng \texttt{product\_variants} với các thuộc tính khác nhau như RAM, dung lượng, màu sắc. Bảng variant lưu giá bán (\texttt{price}), số lượng tồn kho và mã \texttt{sku} duy nhất; giá gốc (\texttt{base\_price}) được lưu ở bảng \texttt{products} vì là đặc tính của sản phẩm, không phải của từng biến thể. Thiết kế này tách biệt thông tin sản phẩm chung (tên, ảnh, mô tả) và thông tin biến thể (giá, kho), cho phép một sản phẩm có nhiều mức giá khác nhau.

10 bảng cốt lõi sử dụng khóa chính INTEGER auto-increment. Thiết kế đáng chú ý: \texttt{orders} nhúng trực tiếp thông tin thanh toán (tránh JOIN), \texttt{order\_items} snapshot giá tại thời điểm mua, \texttt{carts} hỗ trợ cả guest (qua \texttt{session\_id}) và user, \texttt{chat\_messages} lưu lịch sử chatbot kèm intent và thời gian xử lý.

\subsection{Thiết kế hệ thống sản phẩm và biến thể}
\label{sec:c3_product_variant}

Hệ thống sản phẩm theo mô hình \textbf{Product--Variant--Attribute} ba tầng: \texttt{products} (thông tin chung), \texttt{product\_variants} (giá, kho, SKU riêng), và hệ thống thuộc tính (\texttt{attribute\_groups} + \texttt{attribute\_values}) cho phép gán bộ thuộc tính riêng từng sản phẩm. Cơ chế sinh tên biến thể tự động từ \texttt{baseName} + các giá trị có \texttt{affectsName=true} đảm bảo đặt tên nhất quán.

\subsection{Bảng phục vụ AI Chatbot và ghi log}

Hình~\ref{fig:erd_ai_log} trình bày sơ đồ ERD của bốn bảng phục vụ AI Chatbot và hệ thống ghi log.

\begin{figure}[H]
  \centering
  \fitfig{figures/c3/erd_ai_log.pdf}
  \caption{Sơ đồ ERD các bảng AI Chatbot và ghi log hệ thống}
  \label{fig:erd_ai_log}
\end{figure}

Bảng \texttt{chat\_messages} lưu từng tin nhắn chatbot kèm \texttt{session\_id}, \texttt{role}, \texttt{intent}, \texttt{response\_time\_ms} và \texttt{is\_fallback} phục vụ analytics (context phiên lưu trong RAM Map, database chỉ ghi log). Bảng \texttt{search\_histories} và \texttt{recently\_viewed} ghi lịch sử tìm kiếm và xem sản phẩm. Bảng \texttt{inventory\_logs} là immutable audit log mọi thay đổi tồn kho với \texttt{change\_type}, \texttt{change\_amount}, \texttt{previous\_stock} và \texttt{new\_stock}.

\section{Thiết kế sơ đồ trạng thái}
\label{sec:c3_state_diagrams}

Hệ thống TechStore mô hình hóa bốn thực thể nghiệp vụ quan trọng dưới dạng sơ đồ trạng thái để kiểm soát chặt chẽ các chuyển trạng thái hợp lệ và ngăn các lỗi logic do thay đổi trạng thái tùy tiện.

\subsection{Trạng thái đơn hàng}
\label{sec:c3_order_states}

Vòng đời của đơn hàng trong TechStore được mô hình hóa như một máy trạng thái hữu hạn (FSM -- Finite State Machine)~\cite{fowler2003patterns} với 5 trạng thái và các chuyển trạng thái được kiểm soát chặt chẽ. Hình~\ref{fig:order_states} thể hiện biểu đồ trạng thái đơn hàng.

\clearpage
\begin{figure}[H]
  \centering
  \fitfig{figures/c3/order_states.pdf}
  \caption{Biểu đồ trạng thái vòng đời đơn hàng TechStore}
  \label{fig:order_states}
\end{figure}

Đơn hàng đi qua 5 trạng thái: \texttt{pending} → \texttt{processing} (IPN thành công) → \texttt{shipped} → \texttt{delivered}. Hủy (\texttt{cancelled}) có thể từ 3 trạng thái đầu, kèm hoàn tồn kho trong transaction. Mọi chuyển trạng thái đều có validation ngăn chuyển bất hợp lệ.

\subsection{Trạng thái thanh toán}
\label{sec:c3_payment_states}

Trạng thái thanh toán tách biệt qua trường \texttt{payment\_status} với bốn giá trị: \texttt{pending}, \texttt{paid}, \texttt{failed}, \texttt{refunded} (Hình~\ref{fig:payment_states}). Thiết kế tách hai trường trạng thái phản ánh đúng thực tế: đơn hàng có thể đang \texttt{processing} trong khi thanh toán đã \texttt{paid}. Cơ chế idempotency ngăn xử lý trùng lặp khi cổng gửi lại IPN nhiều lần.

\begin{figure}[H]
  \centering
  \fitfig{figures/c3/payment_states.pdf}
  \caption{Biểu đồ trạng thái thanh toán TechStore}
  \label{fig:payment_states}
\end{figure}

\subsection{Trạng thái sản phẩm và tài khoản người dùng}
\label{sec:c3_product_user_states}

Ngoài đơn hàng và thanh toán, hệ thống mô hình hóa thêm hai thực thể đơn giản hơn. Sản phẩm có ba trạng thái (\texttt{active}, \texttt{inactive}, \texttt{draft}) chuyển đổi tự do qua lại; chỉ sản phẩm \texttt{active} được phép đặt hàng, xóa sản phẩm là soft delete. Tài khoản người dùng có bốn trạng thái (\texttt{unverified}, \texttt{active}, \texttt{disabled}, \texttt{deleted}); đăng ký email bắt đầu ở \texttt{unverified} (cần OTP), đăng nhập Google OAuth vào thẳng \texttt{active}; trạng thái \texttt{disabled} chặn cả đăng nhập lẫn làm mới token.

\section{Thiết kế các luồng xử lý chính}
\label{sec:c3_sequences}

Phần này trình bày ba luồng xử lý chính: xác thực người dùng, giỏ hàng và đặt hàng thanh toán.

\subsection{Luồng xác thực người dùng}

Hệ thống hỗ trợ hai phương thức xác thực với chiến lược dual-token: access token (JWT, 7 ngày) trong body response và refresh token (30 ngày) trong \texttt{httpOnly} cookie chống XSS. Đăng ký qua email yêu cầu xác thực OTP 6 số (sinh bằng \texttt{crypto.randomInt}, so sánh bằng \texttt{timingSafeEqual} chống timing attack, thông báo lỗi chung chống user enumeration). Google OAuth cho phép bỏ qua OTP, tạo hoặc merge tài khoản tự động.

\subsection{Luồng giỏ hàng và merge khi đăng nhập}

Giỏ hàng phục vụ cả khách vãng lai (localStorage + session cookie) và khách đã đăng nhập (đồng bộ server). Khi đăng nhập, hook \texttt{useCartMerge} tự động gộp giỏ hàng guest vào tài khoản trong transaction với \texttt{SELECT FOR UPDATE}, cộng dồn số lượng cho items trùng và cập nhật giá hiện tại.

\subsection{Luồng đặt hàng và thanh toán trực tuyến}

Luồng đặt hàng và thanh toán là luồng phức tạp nhất trong hệ thống, liên quan đến Order Service, Payment Service và cổng thanh toán bên ngoài. Luồng được chia thành ba giai đoạn: tạo đơn hàng, sinh URL thanh toán, và xử lý kết quả từ cổng.

\clearpage
Hình~\ref{fig:seq_checkout_create_order} mô tả giai đoạn tạo đơn hàng từ khi người dùng xác nhận giỏ hàng.

\begin{figure}[H]
  \centering
  \fitfig{figures/c3/seq_checkout_create_order.pdf}
  \caption{Biểu đồ tuần tự luồng tạo đơn hàng}
  \label{fig:seq_checkout_create_order}
\end{figure}

\begin{sloppypar}
Order Service mở transaction và thực hiện đầu tiên trong transaction đó là hủy các đơn hàng \texttt{pending} cũ của người dùng kèm hoàn tồn kho tương ứng -- bước này nằm trong transaction để đảm bảo tính nguyên tử, tránh tình huống hủy thành công nhưng transaction chính rollback. Tiếp theo, hệ thống thực hiện \texttt{SELECT FOR UPDATE} cho từng biến thể sản phẩm trong giỏ hàng để khóa hàng, kiểm tra số lượng còn đủ, trừ tồn kho và ghi \texttt{inventory\_logs}. Nếu bất kỳ sản phẩm nào hết hàng hoặc ngừng bán, toàn bộ transaction bị rollback ngay lập tức.

Nếu có mã giảm giá, hệ thống cũng khóa bản ghi mã bằng \texttt{SELECT FOR UPDATE} để kiểm tra điều kiện (còn lượt, chưa hết hạn, đủ giá trị đơn tối thiểu) trước khi áp dụng. Cuối cùng, đơn hàng được tạo với \texttt{status=pending} và \linebreak[0]\texttt{paymentStatus=pending}, kèm \texttt{order\_items} snapshot giá tại thời điểm đặt. Với thanh toán thủ công (COD, chuyển khoản, trả góp), giỏ hàng xóa ngay sau COMMIT và email xác nhận được gửi bất đồng bộ. Với thanh toán online, giỏ hàng và bộ đếm mã giảm giá được giữ lại chờ xác nhận từ IPN callback.
\end{sloppypar}

Sau khi tạo đơn, frontend gọi API sinh URL thanh toán (MoMo dùng HMAC-SHA256, VNPay dùng HMAC-SHA512). Kết quả trả về qua IPN callback (server-to-server) và Return URL. Payment Service xác thực chữ ký HMAC, kiểm tra idempotency (dựa trên \texttt{paymentTransactionId}), và cập nhật trạng thái trong transaction. Sau COMMIT, các tác vụ hậu thanh toán (xóa giỏ hàng, tăng bộ đếm mã giảm giá, gửi email) chạy fire-and-forget.

\section{Thiết kế pipeline RAG cho chatbot AI}
\label{sec:c3_rag_pipeline}

\subsection{Tổng quan thiết kế và lý do chọn kiến trúc}

Chatbot AI là thành phần phân biệt TechStore với các hệ thống thương mại điện tử thông thường. Yêu cầu đặt ra là chatbot phải trả lời câu hỏi tiếng Việt tự nhiên về sản phẩm, kể cả câu hỏi viết tắt, sai chính tả hoặc mơ hồ, đồng thời chỉ sử dụng thông tin sản phẩm thực tế trong hệ thống chứ không suy diễn từ kiến thức chung của LLM.

Kiến trúc thuần LLM (prompt-based) không đáp ứng được yêu cầu này vì LLM không biết danh mục sản phẩm thực tế. Kiến trúc RAG cơ bản (chỉ tìm kiếm ngữ nghĩa) có thể bỏ sót sản phẩm khi người dùng gõ tên model cụ thể (``iPhone 17 Pro''). Do đó, hệ thống được thiết kế theo kiến trúc \textit{Advanced RAG} kết hợp tìm kiếm ngữ nghĩa (dense retrieval) và tìm kiếm từ khóa (sparse retrieval) để tận dụng ưu điểm của cả hai phương pháp, đồng thời bổ sung cơ chế tiền xử lý query và rewrite bằng LLM trước khi truy xuất.

Về lưu trữ vector, hệ thống chọn file JSON (\texttt{vector-db.json}) thay vì cơ sở dữ liệu vector chuyên dụng như Pinecone, Weaviate hay pgvector. Với catalog dưới 10.000 sản phẩm, toàn bộ tập vector (ước tính $10.000 \times 1.024$ chiều $\times$ 4 byte $\approx$ 40MB) có thể nạp vào RAM một lần và tính cosine similarity bằng phép nhân ma trận, đạt độ trễ dưới 50ms mà không cần thêm dịch vụ bên ngoài. Nhược điểm là tìm kiếm theo thứ tự $O(n)$ tuyến tính vì không có chỉ mục HNSW; khi catalog vượt ngưỡng 100.000 sản phẩm hoặc cần độ trễ dưới 10ms, hệ thống có thể chuyển sang pgvector hoặc Qdrant mà chỉ cần thay thế lớp \texttt{HybridVectorStore} mà không ảnh hưởng đến phần còn lại của pipeline.

\subsection{Sơ đồ luồng xử lý pipeline}

Hình~\ref{fig:rag_pipeline_flow} trình bày toàn bộ luồng xử lý 7 bước của pipeline RAG TechStore, từ xác thực đầu vào đến lưu trữ lịch sử hội thoại.

\begin{figure}[H]
  \centering
  \fitfig[0.95]{figures/c3/rag_pipeline_flow.pdf}
  \caption{Sơ đồ luồng xử lý pipeline RAG TechStore (7 bước)}
  \label{fig:rag_pipeline_flow}
\end{figure}

\subsection{Biểu đồ tuần tự pipeline RAG}
\label{sec:c3_rag_sequence}

Pipeline RAG được chia thành ba giai đoạn: tiền xử lý (Hình~\ref{fig:seq_chatbot_preprocess}), truy xuất (Hình~\ref{fig:seq_chatbot_retrieve}), và sinh phản hồi (Hình~\ref{fig:seq_chatbot_generate}).

\begin{figure}[H]
  \centering
  \fitfig{figures/c3/seq_chatbot_preprocess.pdf}
  \caption{Biểu đồ tuần tự chatbot RAG -- Bước 1--3: Tiền xử lý, phân loại và kiểm tra bảo mật}
  \label{fig:seq_chatbot_preprocess}
\end{figure}

\begin{figure}[H]
  \centering
  \fitfig{figures/c3/seq_chatbot_retrieve.pdf}
  \caption{Biểu đồ tuần tự chatbot RAG -- Bước 4--5: Load session và Hybrid Search}
  \label{fig:seq_chatbot_retrieve}
\end{figure}

\begin{figure}[H]
  \centering
  \fitfig{figures/c3/seq_chatbot_generate.pdf}
  \caption{Biểu đồ tuần tự chatbot RAG -- Bước 6--7: Sinh phản hồi và lưu trữ}
  \label{fig:seq_chatbot_generate}
\end{figure}

\subsection{Các bước trong pipeline RAG}

Pipeline RAG gồm 7 bước tuần tự với 4 điểm phân nhánh (validate, lọc nội dung, retrieval rỗng, LLM timeout):

\textbf{Bước 1--3 (Tiền xử lý):} Validate đầu vào (không rỗng, $\leq$500 ký tự). Chuẩn hóa 71 mẫu regex viết tắt ($<$1ms) và phân loại 6 intent. Lọc prompt injection (15 nhóm, 24 mẫu regex) và off-topic bằng regex ($<$10ms), kết thúc sớm không gọi LLM.

\textbf{Bước 4 (Session):} Lấy lịch sử từ in-memory Map (tối đa 10 lượt, TTL 30 phút, LRU 500 phiên). Nếu câu hỏi chứa đại từ (``cái đó'', ``nó''), bổ sung tên sản phẩm từ tin nhắn trước vào query.

\textbf{Bước 5 (Retrieval):} \texttt{Promise.all} chạy đồng thời \texttt{hybridSearch(query, topK=10, minScore=0.45)} và \texttt{rewriteQuery} (LLM, timeout 8s). Nếu rỗng, fallback \texttt{hybridSearch(query, 3, 0)} với \texttt{lowConfidence=true}.

\textbf{Bước 6 (Generation):} \texttt{buildAugmentedPrompt} ghép sản phẩm + lịch sử + system prompt gửi LLM (\texttt{temperature=0.3}, JSON output). \texttt{parseLLMOutput} khớp sản phẩm qua 4 mức (exact → version → model number → 80\% từ). Timeout hoặc lỗi → \texttt{simpleKeywordMatch} làm fallback.

\textbf{Bước 7 (Persist):} Fire-and-forget cập nhật session Map và \texttt{ChatMessage.bulkCreate} vào MySQL.

\subsection{Thiết kế Hybrid Search}
\label{sec:c3_hybrid_search}

\texttt{hybridSearch(query, topK, minScore)} thực hiện ba bước: (1) embedding query thành vector 1024 chiều qua chain fallback, tính cosine similarity với toàn bộ product vectors (chỉ giữ $\geq$ minScore); (2) tìm kiếm từ khóa BM25-inspired với trọng số tên sản phẩm $\times$3; (3) hợp nhất kết quả, overlap boost +0.05 cho sản phẩm xuất hiện ở cả hai, đánh dấu \texttt{lowConfidence} cho kết quả chỉ từ keyword. Pipeline gọi ở hai chế độ: chính (topK=10, minScore=0.45) và dự phòng (topK=3, minScore=0).

\subsection{Session memory, vector sync và embedding fallback}

Session memory lưu trong \texttt{Map} (tối đa 500 phiên, TTL 30 phút, LRU eviction) — phù hợp quy mô hiện tại nhưng mất khi server restart. Vector store tự động rebuild khi chênh lệch $>$5\% so với catalog, và cập nhật tăng dần qua Sequelize model hooks (\texttt{afterCreate/Update/Destroy}).

Embedding chain fallback qua ba provider đều xuất vector 1024 chiều (Hình~\ref{fig:embedding_fallback}): Jina v3 (ưu tiên, hỗ trợ task type riêng cho query/passage) → HuggingFace e5-large-instruct → e5-large. Nếu tất cả thất bại, pipeline fallback về \texttt{simpleKeywordMatch}.

\begin{figure}[H]
  \centering
  \fitfig[0.90]{figures/c3/embedding_fallback.pdf}
  \caption{Sơ đồ luồng chuỗi fallback embedding với ba nhà cung cấp}
  \label{fig:embedding_fallback}
\end{figure}

\bigskip
Chương này đã trình bày thiết kế hệ thống TechStore từ phân tích yêu cầu, kiến trúc Modular Monolith và Feature-Based, cơ sở dữ liệu, sơ đồ trạng thái, các luồng xử lý chính, đến pipeline RAG Hybrid 7 bước với Hybrid Search và embedding chain fallback. Chương~4 sẽ trình bày kết quả cài đặt, kiểm thử và đánh giá hiệu quả thực tế của hệ thống.
